\documentclass[11pt,a4paper]{article}

\usepackage{kotex}
\usepackage{fontspec}
\setmainfont{Times New Roman}
\setmainhangulfont{AppleMyungjo}
\setsanshangulfont{Apple SD Gothic Neo}

\usepackage[a4paper,top=16mm,bottom=15mm,left=20mm,right=20mm]{geometry}
\usepackage{fancyhdr}
\setlength{\headheight}{24pt}
\usepackage{enumitem}
\usepackage{mdframed}
\usepackage{array}
\usepackage{booktabs}

\setlength{\parindent}{1em}
\linespread{1.08}

\pagestyle{fancy}
\fancyhf{}
\renewcommand{\headrulewidth}{0.6pt}
\fancyhead[L]{\sffamily\bfseries 2026 고1 영어 변형 --- 지문 43-45}
\fancyhead[R]{\sffamily [Emma와 Lisa]}
\renewcommand{\footrulewidth}{0pt}
\fancyfoot[C]{\framebox[16mm][c]{\thepage}}

\newcommand{\qno}[1]{\par\vspace{8pt}\noindent{\sffamily\bfseries #1.}\ }
\newcommand{\instr}[1]{{\sffamily #1}\par\vspace{3pt}}
\newlist{choices}{enumerate}{1}
\setlist[choices]{label=,leftmargin=1.5em,itemsep=2pt,topsep=3pt,parsep=0pt}
\newmdenv[linewidth=0.6pt,roundcorner=3pt,innertopmargin=7pt,innerbottommargin=7pt,
          innerleftmargin=9pt,innerrightmargin=9pt,skipabove=5pt,skipbelow=5pt]{pbox}
\newcommand{\nemo}[1]{\fbox{\,#1\,}}
\newcommand{\secl}[1]{\par\vspace{5pt}\noindent{\sffamily\bfseries #1}\ }

\begin{document}

\begin{center}
{\fontsize{15}{17}\selectfont\sffamily\bfseries [지문 43-45] 다음 글을 읽고, 물음에 답하시오.}
\end{center}
\vspace{4pt}

\begin{pbox}
\noindent{\sffamily\bfseries (A)}\ \
Emma was a talented vocalist but struggled to gain recognition. Lisa, who went to the
same high school with her, had quickly become famous with powerful performances. Emma
felt jealous watching Lisa's success. One day, Emma decided to enter a major vocal
competition in the capital city. \underline{(a) She} spent months preparing her masterpiece
of a performance, hoping this would finally be her opportunity to shine.
\end{pbox}

\secl{(B)}
Moments later, Lisa came over to congratulate Emma with a warm smile. She said, ``Your
hard work and dedication really paid off. I'm so happy for you.'' \underline{(b) She} didn't
seem bothered by the results. ``I'm proud that your singing brings joy to people,'' said
Lisa. Emma felt grateful for Lisa's kindness. Lisa added, ``I've always been moved by your
effort and have been supporting \underline{(c) you} all along.''

\secl{(C)}
When the competition day arrived, Emma felt nervous. Walking into the concert hall, Emma
saw a large audience cheering for Lisa. She couldn't help but compare herself to Lisa. On
her turn, Emma performed with all her heart, pouring \underline{(d) her} months of
preparation into every note. The judges then evaluated each vocalist's performance
carefully. Finally, the results were announced. To her surprise, Emma got first place, not
Lisa.

\secl{(D)}
These words made Emma feel embarrassed at her jealousy. At the same time, \underline{(e) she}
got inspired to find a better path forward. They hugged warmly, their bond strengthened.
From that day on, they supported each other and even sang on the same stage together. Emma
became more confident and, finally, got a chance to have her solo concert.

\vspace{4pt}

\qno{1}\instr{주어진 글 (A)에 이어질 내용을 순서에 맞게 배열한 것으로 가장 적절한 것은? [2점]}
\begin{choices}
\item ① (B)-(D)-(C)
\item ② (C)-(B)-(D)
\item ③ (C)-(D)-(B)
\item ④ (D)-(B)-(C)
\item ⑤ (D)-(C)-(B)
\end{choices}

\vspace{4pt}

\qno{2}\instr{밑줄 친 (a)$\sim$(e) 중에서 가리키는 대상이 나머지 넷과 다른 것은? [2점]}
\begin{choices}
\item ① (a) \quad ② (b) \quad ③ (c) \quad ④ (d) \quad ⑤ (e)
\end{choices}

\vspace{4pt}

\qno{3}\instr{윗글의 내용과 일치하지 \emph{않는} 것은? [2점]}
\begin{choices}
\item ① Both young women received their education at the same secondary institution.
\item ② The more renowned of the two singers approached the other to offer kind words after
the competition.
\item ③ Upon entering the performance venue, Emma observed that a substantial number of
audience members were showing enthusiasm for her rival.
\item ④ Emma secured the runner-up position in the competition, missing out on the top prize.
\item ⑤ Following the competition, the two vocalists had an opportunity to share the same stage.
\end{choices}

\vspace{4pt}

\qno{4}\instr{[서술형] Emma의 심경이 대회 전후로 어떻게 변화했는지 우리말로 한 문장으로 쓰시오.
(심경을 나타내는 단어 두 개 이상 사용)}
\par\vspace{2pt}
\noindent$\Rightarrow$ \rule{0.89\linewidth}{0.4pt}

\end{document}
